\clearpage\pagestyle{otMain}\fontsize{11}{15}\selectfont
\corrigesChapitre{9}{Repérage et configurations}
\begin{blocLivre}{cor-9-1}
\begin{corrige}{1}{Distance entre deux points}
\demarche{Calculer le carré de la distance, puis prendre la racine positive.}
Les différences de coordonnées sont 4 et 3. Dans le repère orthonormé, $AB^2=4^2+3^2=25$, donc $\res{AB=5}$. La racine positive est retenue puisqu’une longueur est positive. Le triangle auxiliaire du dessin relie ces différences au théorème de Pythagore.
\label{fig-cor-9-1}
\figureGeo[corrige]{conf-01}
\end{corrige}
\end{blocLivre}
\begin{blocLivre}{cor-9-2}
\begin{corrige}{2}{Coordonnées d’un milieu}
\demarche{Faire deux moyennes et contrôler le placement sur le segment.}
On fait la moyenne des coordonnées : $I((-2+6)/2;(5-1)/2)=\res{I(2;2)}$. Puis $\vect{AI}=(4;-3)$ et $\vect{IB}=(4;-3)$. Cette égalité confirme que $I$ est bien le milieu, et pas seulement un point équidistant des deux extrémités.
\label{fig-cor-9-2}
\figureGeo[corrige]{conf-02}
\end{corrige}
\end{blocLivre}
\begin{blocLivre}{cor-9-3}
\begin{corrige}{3}{Retrouver une extrémité}
\demarche{Résoudre les deux formules du milieu pour l’extrémité inconnue.}
Les formules de milieu donnent $(-1+x_B)/2=2$ et $(4+y_B)/2=1$. Donc $x_B=5$ et $y_B=-2$ : $\res{B(5;-2)}$. Vérification : la moyenne des abscisses vaut 2 et celle des ordonnées 1. Le point $B$ est le symétrique de $A$ par rapport à $I$.
\label{fig-cor-9-3}
\figureGeo[corrige]{conf-03}
\end{corrige}
\end{blocLivre}
\begin{blocLivre}{cor-9-4}
\begin{corrige}{4}{Un triangle rectangle}
\demarche{Vérifier Pythagore puis utiliser les deux côtés perpendiculaires.}
$AB^2=16$, $BC^2=9$ et $AC^2=25$. Puisque $AB^2+BC^2=AC^2$, la réciproque de Pythagore montre que le triangle est rectangle en $B$. L’aire vaut $AB\times BC/2=4\times3/2=\res6$. Le côté $AC$, opposé à l’angle droit, est bien le plus long.
\label{fig-cor-9-4}
\figureGeo[corrige]{conf-04}
\end{corrige}
\end{blocLivre}
\begin{blocLivre}{cor-9-5}
\begin{corrige}{5}{Une médiatrice facile}
\demarche{Associer une médiatrice à un milieu et à une perpendicularité.}
Le milieu est $(1;0)$. Le segment étant horizontal, sa médiatrice est verticale : $\res{x=1}$. Le point $M(1;3)$ lui appartient. On contrôle l’équidistance : $MA^2=3^2+3^2=18$ et $MB^2=(-3)^2+3^2=18$.
\label{fig-cor-9-5}
\figureGeo[corrige]{conf-05}
\end{corrige}
\end{blocLivre}
\begin{blocLivre}{cor-9-6}
\begin{corrige}{6}{Un triangle isocèle}
\demarche{Comparer les carrés des côtés et lire une hauteur perpendiculaire.}
$CA^2=2^2+3^2=13$ et $CB^2=(-2)^2+3^2=13$. Ainsi $CA=CB=\sqrt{13}$ : le triangle est isocèle en $C$. Son périmètre est $\res{4+2\sqrt{13}}$. La hauteur sur la base horizontale vaut 3, donc l’aire est $4\times3/2=\res6$.
\label{fig-cor-9-6}
\figureGeo[corrige]{conf-06}
\end{corrige}
\end{blocLivre}
\begin{blocLivre}{cor-9-7}
\begin{corrige}{7}{Thalès dans un triangle}
\demarche{Identifier le coefficient de réduction et le mettre au carré pour les aires.}
Le parallélisme autorise Thalès : $AM/AB=AN/AC=2/3$. Ainsi $AN=12\times2/3=\res8$. Les longueurs sont multipliées par $2/3$, donc les aires par son carré : $\res{\mathcal A(AMN)/\mathcal A(ABC)=4/9}$. Le facteur d’aire n’est pas le facteur de longueur.
\label{fig-cor-9-7}
\figureGeo[corrige]{conf-07}
\end{corrige}
\end{blocLivre}
\begin{blocLivre}{cor-9-8}
\begin{corrige}{8}{Une hauteur par trigonométrie}
\demarche{Choisir le rapport trigonométrique correspondant aux côtés connus.}
Le triangle est rectangle en $B$. Pour l’angle en $P$, la hauteur est le côté opposé et 20 le côté adjacent : $\tan35^\circ=BS/20$. Donc $\res{BS=20\tan35^\circ\approx14{,}0\text{ m}}$. La calculatrice doit être réglée en degrés ; il faut ensuite arrondir, sans modifier les données en cours de calcul.
\label{fig-cor-9-8}
\figureGeo[corrige]{conf-08}
\end{corrige}
\end{blocLivre}
\begin{blocLivre}{cor-9-9}
\begin{corrige}{9}{Une réduction de cône}
\demarche{Utiliser la similitude puis la formule adaptée à chaque solide.}
Le petit cône est une réduction de coefficient $1/2$ : son rayon vaut 1,5 cm et sa hauteur 4 cm. Son volume est $\pi\times1{,}5^2\times4/3=\res{3\pi\text{ cm}^3}$. Le grand volume vaut $24\pi$ cm$^3$. Le rapport est $\res{1/8}$, soit le cube du coefficient de réduction. Pour la boule, $4\pi\times3^3/3=\res{36\pi\text{ cm}^3}$. Pour la pyramide, $27\times8/3=\res{72\text{ cm}^3}$. Les trois formules concernent des solides différents ; l’aire de base seule ne détermine pas le volume.
\label{fig-cor-9-9}
\figureGeo[corrige]{conf-09}
\end{corrige}
\end{blocLivre}
\begin{blocLivre}{cor-9-10}
\begin{corrige}{10}{Le centre d’un parallélogramme}
\demarche{Calculer les milieux des deux diagonales et appliquer la réciproque.}
Le milieu de $[AC]$ est $(5/2;2)$. Celui de $[BD]$ est $((4+1)/2;(1+3)/2)=(5/2;2)$. Les diagonales ont le même milieu : le quadrilatère $ABCD$ est donc un parallélogramme. Le critère concerne les diagonales $AC$ et $BD$, pas deux côtés quelconques.
\label{fig-cor-9-10}
\figureGeo[corrige]{conf-10}
\end{corrige}
\end{blocLivre}
\begin{blocLivre}{cor-9-11}
\begin{corrige}{11}{Un rectangle incliné}
\demarche{Combiner un critère de parallélogramme et la réciproque de Pythagore.}
On a $\vect{AB}=(4;2)=\vect{DC}$, donc c’est un parallélogramme. Puis $AB^2=20$, $BC^2=5$ et $AC^2=25$. L’égalité $AB^2+BC^2=AC^2$ prouve que l’angle en $B$ est droit. Un parallélogramme avec un angle droit est un rectangle. L’orientation des côtés sur la page n’intervient pas.
\label{fig-cor-9-11}
\figureGeo[corrige]{conf-11}
\end{corrige}
\end{blocLivre}
\begin{blocLivre}{cor-9-12}
\begin{corrige}{12}{Un triangle avec un paramètre}
\demarche{Traduire l’équidistance par une égalité de distances au carré.}
On compare les carrés : $(t+1)^2+4=(t-3)^2+4$. Après développement, $2t+1=-6t+9$, soit $t=1$. Le point $C(1;2)$ se trouve sur la médiatrice verticale de $[AB]$. La base vaut 4 et la hauteur 2, donc l’aire est $\res4$.
\label{fig-cor-9-12}
\figureGeo[corrige]{conf-12}
\end{corrige}
\end{blocLivre}
\begin{blocLivre}{cor-9-13}
\begin{corrige}{13}{Un point équidistant}
\demarche{Intersecter deux médiatrices simples puis contrôler la troisième distance.}
La médiatrice de $[AB]$ est $x=2$ et celle de $[AC]$ est $y=3$. Leur intersection est $\res{O(2;3)}$. On calcule $OA^2=13$, $OB^2=13$ et $OC^2=13$. Le cercle de centre $O$ et de rayon $\res{\sqrt{13}}$ passe donc par les trois points. Les égalités de distance vérifient directement la construction.
\label{fig-cor-9-13}
\figureGeo[corrige]{conf-13}
\end{corrige}
\end{blocLivre}
\begin{blocLivre}{cor-9-14}
\begin{corrige}{14}{Un lieu de points sur une horizontale}
\demarche{Résoudre la distance au carré puis distinguer longueur du cercle et aire du disque.}
On écrit $OM^2=t^2+16=25$, donc $t^2=9$ et $t=-3$ ou 3. Les points sont $\res{(-3;4)\text{ et }(3;4)}$. Ils sont symétriques par rapport à l’axe des ordonnées. Garder seulement la racine positive aurait oublié une intersection. Le cercle a pour périmètre $2\pi\times5=\res{10\pi}$ ; le disque a pour aire $\pi\times5^2=\res{25\pi}$. Une unité de longueur devient une unité d’aire dans le dernier résultat.
\label{fig-cor-9-14}
\figureGeo[corrige]{conf-14}
\end{corrige}
\end{blocLivre}
\begin{blocLivre}{cor-9-15}
\begin{corrige}{15}{Un raccourci dans un pavé}
\demarche{Ramener une question spatiale à deux triangles rectangles successifs.}
La diagonale de base vaut $\sqrt{6^2+8^2}=10$ cm. Elle forme avec la hauteur un triangle rectangle, donc la diagonale intérieure vaut $\sqrt{10^2+5^2}=\res{5\sqrt5\text{ cm}}$, environ 11,18 cm. Le trajet de trois arêtes vaut $6+8+5=19$ cm ; il est plus long. Les deux applications de Pythagore sont faites dans deux plans perpendiculaires appropriés.
\label{fig-cor-9-15}
\figureGeo[corrige]{conf-15}
\end{corrige}
\end{blocLivre}
